\documentclass[aspectratio=169,11pt]{beamer}

\usetheme{Madrid}
\usecolortheme{default}
\usepackage{amsmath,amssymb}
\usepackage{graphicx}
\usepackage{booktabs}
\usepackage{tikz}
\usepackage{pgfplots}
\pgfplotsset{compat=1.18}

\definecolor{classicalblue}{RGB}{33,150,243}
\definecolor{dlred}{RGB}{255,87,34}

\title{Benchmarking ML/DL Models for\\Lithium-Ion Battery SOH Estimation}
\author{Malak Mekyassi}
\institute{Capgemini Engineering}
\date{August 2025}

\begin{document}

\frame{\titlepage}

\begin{frame}{Outline}
\tableofcontents
\end{frame}

\section{Problem \& Motivation}

\begin{frame}{Why SOH Estimation Matters}
\begin{itemize}[nosep]
    \item \textbf{Safety:} detect abnormal degradation before thermal events
    \item \textbf{Warranty:} fair residual value for second-life reuse
    \item \textbf{Charging:} adaptive strategies to extend battery lifetime
    \item \textbf{Range:} accurate remaining capacity estimates for drivers
\end{itemize}
\vspace{0.5em}
\textbf{Challenge:} SOH cannot be measured directly in deployed vehicles. BMS only has access to voltage, current, and temperature.
\end{frame}

\begin{frame}{Research Question}
\begin{block}{How do classical ML and DL models compare for SOH estimation?}
\begin{itemize}[nosep]
    \item 4 classical models: Naive, RF, SVR, GPR
    \item 3 DL models: LSTM, 1D CNN, Transformer
    \item 2 open datasets: NASA PCoE and CALCE CS2
    \item Cell-based LOOCV to prevent data leakage
\end{itemize}
\end{block}
\end{frame}

\section{Datasets}

\begin{frame}{Datasets Overview}
\begin{table}[H]
\centering
\small
\begin{tabular}{lcc}
\toprule
\textbf{Property} & \textbf{NASA PCoE} & \textbf{CALCE CS2} \\
\midrule
Cell chemistry & LCO (prismatic) & LCO (coin) \\
Cells & 4 (B0005--B0018) & 4 (CS2\_33--CS2\_36) \\
Rated capacity & 2.0 Ah & 1.1 Ah \\
Discharge cycles/cell & 168 & 775--973 \\
Total cycles & 636 & 3,543 \\
EIS measurements & Yes & No \\
\bottomrule
\end{tabular}
\end{table}
\vspace{0.5em}
\textbf{Key difference:} CALCE has $\sim$6$\times$ more training data than NASA.
\end{frame}

\section{Methodology}

\begin{frame}{Preprocessing Pipeline}
\begin{enumerate}[nosep]
    \item \textbf{Cycle segmentation} --- separate charge/discharge/impedance
    \item \textbf{Capacity computation} --- current integration via trapezoidal rule
    \item \textbf{Voltage filtering} --- Savitzky-Golay (window=51, order=3)
    \item \textbf{Resampling} --- uniform capacity grid (1,000 points)
    \item \textbf{SOH labeling} --- $SOH(n) = Q_{discharge}(n) / Q_{initial}$
    \item \textbf{Partial cycle filtering} --- discard early cycles $<$ 90\% of Q\_initial
\end{enumerate}
\end{frame}

\begin{frame}{Feature Engineering --- 12 Features}
\begin{columns}
\column{0.5\textwidth}
\textbf{Physical Features}
\begin{itemize}[nosep,small]
    \item ICA: peak voltage, height, area, secondary ratio
    \item Internal resistance: $\Delta V / \Delta I$
    \item Discharge energy: $\int V \cdot I \, dt$
    \item Mean discharge voltage
    \item Coulombic efficiency: $Q_{dis} / Q_{chg}$
    \item Temperature: mean, range
    \item Capacity fade rate
\end{itemize}

\column{0.5\textwidth}
\textbf{Feature Selection}
\begin{enumerate}[nosep,small]
    \item Correlation filter ($|r| > 0.95$)
    \item RF importance ranking (top 10)
\end{enumerate}
\vspace{1em}
Top features:
\begin{enumerate}[nosep,small]
    \item discharge\_energy
    \item mean\_discharge\_voltage
    \item internal\_resistance
    \item capacity\_fade\_rate
\end{enumerate}
\end{columns}
\end{frame}

\begin{frame}{Models \& Validation}
\begin{columns}
\column{0.5\textwidth}
\textbf{Classical ML}
\begin{itemize}[nosep,small]
    \item \textbf{RF:} 200 trees, depth 20
    \item \textbf{SVR:} RBF, $C$=100, $\epsilon$=0.01
    \item \textbf{GPR:} Mat\'ern 1/2 kernel
    \item \textbf{Naive:} predict mean
\end{itemize}

\column{0.5\textwidth}
\textbf{Deep Learning}
\begin{itemize}[nosep,small]
    \item \textbf{LSTM:} 2-layer (64$\to$32)
    \item \textbf{1D CNN:} 3 conv layers
    \item \textbf{Transformer:} 2 blocks, 4 heads
\end{itemize}
\vspace{0.5em}
\textbf{Validation:} Cell-based LOOCV\\
\textbf{Window:} 20 consecutive cycles
\end{columns}
\end{frame}

\section{Results}

\begin{frame}{CALCE CS2 Results}
\begin{table}[H]
\centering
\small
\begin{tabular}{lcccc}
\toprule
\textbf{Model} & \textbf{RMSE} & \textbf{MAE} & \textbf{MaxAE} & \textbf{$R^2$} \\
\midrule
Naive & 0.1761 & 0.1236 & 0.6234 & $-$0.147 \\
\textbf{RF} & \textbf{0.0226} & \textbf{0.0137} & \textbf{0.0905} & \textbf{0.969} \\
SVR & 0.0593 & 0.0258 & 0.4388 & 0.870 \\
GPR & 0.0368 & 0.0151 & 0.4386 & 0.844 \\
\midrule
LSTM & 0.0611 & --- & --- & 0.781 \\
1D CNN & 0.1438 & --- & --- & $-$0.062 \\
\textbf{Transformer} & \textbf{0.0425} & --- & --- & \textbf{0.914} \\
\bottomrule
\end{tabular}
\end{table}
\vspace{-0.5em}
\begin{itemize}[nosep,small]
    \item RF best overall (RMSE~=~0.023, $R^2$~=~0.969)
    \item Transformer competitive with sufficient data (RMSE~=~0.043, $R^2$~=~0.914)
\end{itemize}
\end{frame}

\begin{frame}{NASA PCoE Results}
\begin{table}[H]
\centering
\small
\begin{tabular}{lcccc}
\toprule
\textbf{Model} & \textbf{RMSE} & \textbf{MAE} & \textbf{MaxAE} & \textbf{$R^2$} \\
\midrule
Naive & 0.1120 & 0.0965 & 0.2029 & $-$0.285 \\
\textbf{RF} & \textbf{0.0266} & 0.0196 & 0.0838 & \textbf{0.918} \\
SVR & 0.0599 & 0.0106 & 0.0540 & 0.580 \\
GPR & 0.0366 & 0.0272 & 0.0977 & 0.835 \\
\midrule
LSTM & 0.0476 & --- & --- & $-$0.135 \\
1D CNN & 0.1040 & --- & --- & $-$2.504 \\
Transformer & 0.0673 & --- & --- & $-$0.142 \\
\bottomrule
\end{tabular}
\end{table}
\vspace{-0.5em}
\begin{itemize}[nosep,small]
    \item RF again best (RMSE~=~0.027, $R^2$~=~0.918)
    \item All DL models have negative $R^2$ --- 636 cycles too few for DL
\end{itemize}
\end{frame}

\begin{frame}{Cross-Dataset Comparison}
\begin{columns}
\column{0.5\textwidth}
\begin{center}
\textbf{NASA PCoE (636 cycles)}\\[0.5em]
\begin{tabular}{cc}
\toprule
Model & $R^2$ \\
\midrule
RF & 0.918 \\
GPR & 0.835 \\
SVR & 0.580 \\
Transformer & $-$0.142 \\
\bottomrule
\end{tabular}
\end{center}

\column{0.5\textwidth}
\begin{center}
\textbf{CALCE CS2 (3,543 cycles)}\\[0.5em]
\begin{tabular}{cc}
\toprule
Model & $R^2$ \\
\midrule
RF & 0.969 \\
Transformer & 0.914 \\
SVR & 0.870 \\
GPR & 0.844 \\
\bottomrule
\end{tabular}
\end{center}
\end{columns}
\vspace{1em}
\textbf{Key insight:} DL needs $\geq$3,000+ cycles to become competitive. Classical ML works well with $<$1,000 cycles.
\end{frame}

\section{Discussion}

\begin{frame}{Key Findings}
\begin{enumerate}[nosep]
    \item \textbf{RF is the best overall model:} RMSE~=~0.023 on CALCE, 0.027 on NASA. Robust, fast, interpretable.
    \item \textbf{DL needs data:} Transformer achieves $R^2$~=~0.914 on CALCE (3,543 cycles) but $R^2$~=~$-$0.142 on NASA (636 cycles).
    \item \textbf{CNN fails:} Insufficient temporal structure in per-cycle feature vectors for convolutional learning.
    \item \textbf{Feature importance:} discharge\_energy, mean\_voltage, internal\_resistance, fade\_rate are most predictive --- aligns with electrochemistry.
    \item \textbf{Deployability:} Classical models: 0.1--0.4s inference. DL: 5--16s. RF/SVR preferred for embedded BMS.
\end{enumerate}
\end{frame}

\begin{frame}{Practical Recommendations}
\begin{block}{For BMS Deployment}
\begin{itemize}[nosep]
    \item \textbf{RF} or \textbf{SVR} with handcrafted features --- best accuracy, fastest inference
    \item \textbf{Transformer} if large training dataset is available ($\geq$3,000 cycles)
    \item \textbf{Avoid CNN} for this task --- consistently underperforms
\end{itemize}
\end{block}
\vspace{1em}
\begin{block}{For Research}
\begin{itemize}[nosep]
    \item Cross-dataset transfer learning as next step
    \item Calendar aging effects not captured in cycling data
    \item More diverse cells/conditions needed for DL generalization
\end{itemize}
\end{block}
\end{frame}

\begin{frame}{Conclusion}
\begin{itemize}[nosep]
    \item Systematic benchmark of 7 models on 2 datasets ($\sim$4,200 cycles total)
    \item Random Forest achieves state-of-the-art on both datasets
    \item Deep learning shows promise but requires $\geq$6$\times$ more data than classical ML
    \item Feature engineering is critical: 12 physical features capture degradation physics
    \item Open-source pipeline for reproducible SOH benchmarking
\end{itemize}
\vspace{1em}
\begin{center}
\textbf{Thank you!}\\
\vspace{0.5em}
\small Code: \texttt{github.com/mal0101/soh\_estimation}
\end{center}
\end{frame}

\end{document}
